\chapter{Mặc Tử \& Hàn Phi — Kiêm Ái, Pháp Trị Và Trật Tự Xã Hội}
\markboth{Mặc Tử \& Hàn Phi — Kiêm Ái, Pháp Trị Và Trật Tự Xã Hội}{Mozi \& Han Fei — Universal Love, Rule of Law, and Social Order}

\section{Kiêm Ái — Yêu Không Phân Biệt}
\begin{truyen}{Kiêm Ái — Yêu Không Phân Biệt}{Jian Ai — Universal Love Without Distinction}
\chuhoa{M}{ặc Tử hỏi học trò: ``Vì sao thiên hạ loạn?''}
\textit{(Mozi asked his students: ``Why is the world in disorder?'')}

Học trò đáp: ``Vì người người không yêu nhau.''
\textit{(Students replied: ``Because people do not love one another.'')}

``Đúng. Và tại sao người người không yêu nhau?'' Mặc Tử tiếp. ``Vì họ yêu người nhà mình hơn người ngoài, yêu nước mình hơn nước khác. Đây là tình yêu phân biệt — và nó tạo ra xung đột.''
\textit{(``Correct. And why do people not love one another?'' Mozi continued. ``Because they love their own family more than others, their own country more than other countries. This is partial love — and it creates conflict.'')}

``Giải pháp là Kiêm Ái — yêu người khác như yêu người nhà mình. Coi nước người khác như nước mình.''
\textit{(``The solution is Jian Ai — loving others as you love your own family. Regarding another's country as your own.'')}

Học trò phản đối: ``Điều đó phi thực tế — người ta không thể yêu người lạ như người thân.''
\textit{(Students objected: ``That is impractical — people cannot love strangers as family.'')}

Mặc Tử đáp: ``Hãy nhìn người lính đi chinh chiến. Ông ta muốn người bạn tốt nhất trông nom gia đình mình. Ai sẽ được chọn? Người chỉ lo gia đình mình hay người yêu thương mọi người như nhau?'' Cả lớp im lặng suy nghĩ.
\textit{(Mozi replied: ``Consider a soldier going to war. He wants the best person to look after his family. Who will he choose? The one who only cares for his own family, or the one who loves all people equally?'' The class fell silent in thought.)}

\end{truyen}

\begin{baihoc}
Kiêm Ái không đòi hỏi ta xóa bỏ sự gắn bó với gia đình — mà mở rộng vòng tròn quan tâm. Sự tiến bộ đạo đức là mở rộng ``chúng ta'' để bao gồm ngày càng nhiều người hơn.
\textit{(Universal Love does not require erasing attachment to family — but expanding the circle of concern. Moral progress is expanding ``us'' to include ever more people.)}
\end{baihoc}
\ngancach

\section{Mặc Tử Phản Chiến — Ý Chí Trời Không Muốn Chiến}
\begin{truyen}{Mặc Tử Phản Chiến — Ý Chí Trời Không Muốn Chiến}{Mozi Against War — Heaven's Will Does Not Want War}
\chuhoa{V}{ua nước Sở chuẩn bị tấn công nước Tống. Mặc Tử nghe tin, đi bộ mười ngày mười đêm đến gặp vua.}
\textit{(The King of Chu was preparing to attack the state of Song. When Mozi heard this, he walked for ten days and nights to see the king.)}

``Đại vương,'' Mặc Tử nói, ``nếu một người bỏ xe ngựa của mình để đi ăn trộm xe đẩy của hàng xóm — người đó là kẻ trộm hay không?''
\textit{(``Your Majesty,'' Mozi said, ``if a man abandons his own carriage to steal his neighbor's cart — is that man a thief?'')}

``Tất nhiên là thẻ,'' vua trả lời.
\textit{(``Of course,'' the king replied.)}

``Nước Sở đất rộng dân đông, nước Tống nhỏ yếu. Đại vương muốn chiếm nước Tống — đó là điều gì nếu không phải trộm cướp?'' Vua không trả lời được.
\textit{(``Chu is vast and populous; Song is small and weak. Your majesty wishes to take Song — what is that if not robbery?'' The king could not answer.)}

Mặc Tử tiếp: ``Thiên hạ không có ác lớn hơn chiến tranh tấn công. Giết một người là tội ác — giết một vạn người lại được gọi là anh hùng? Logic đó sai từ gốc.''
\textit{(Mozi continued: ``There is no greater evil in the world than aggressive war. Killing one person is a crime — killing ten thousand is called heroism? That logic is wrong at its root.'')}

\end{truyen}

\begin{baihoc}
Hãy kiểm tra thang đo đạo đức của bạn để tránh thiên kiến quy mô: điều bạn lên án ở cá nhân — bạn có nhất quán lên án ở quốc gia và tổ chức không?
\textit{(Check your moral scale to avoid the bias of magnitude: what you condemn in individuals — do you consistently condemn in nations and organizations?)}
\end{baihoc}
\ngancach

\section{Hàn Phi — Pháp Trị Và Nghệ Thuật Cai Trị}
\begin{truyen}{Hàn Phi — Pháp Trị Và Nghệ Thuật Cai Trị}{Han Fei — Rule of Law and the Art of Governance}
\chuhoa{H}{àn Phi, học trò của Tuân Tử, đưa ra lập luận gây tranh cãi: ``Đức trị thất bại vì nó phụ thuộc vào sự đức hạnh của người cai trị — điều không thể đảm bảo.''}
\textit{(Han Fei, a student of Xunzi, made a controversial argument: ``Rule by virtue fails because it depends on the virtue of the ruler — something that cannot be guaranteed.'')}

``Nhưng luật pháp rõ ràng, có thưởng có phạt bình đẳng — không phụ thuộc vào đức hạnh cá nhân. Đó là hệ thống có thể vận hành dù người cai trị tốt hay xấu.''
\textit{(``But clear laws, with equal rewards and punishments — these do not depend on individual virtue. That is a system that can function whether the ruler is good or bad.'')}

Ông đưa ra ngụ ngôn: ``Cây thẳng không cần thợ mộc tài năng. Cây cong cần thợ giỏi để uốn. Xã hội phụ thuộc vào bậc thánh nhân là xã hội dễ vỡ khi thánh nhân không còn.''
\textit{(He offered a parable: ``A straight tree needs no skilled carpenter. A crooked tree needs a master to shape it. A society that depends on sages is fragile when the sage is gone.'')}

``Hệ thống tốt sẽ biến người bình thường thành người tốt. Hệ thống xấu sẽ biến người tốt thành kẻ xấu. Hãy xây hệ thống, đừng chờ thánh nhân.''
\textit{(``A good system turns ordinary people into good people. A bad system turns good people into bad. Build the system; do not wait for sages.'')}

Tư tưởng của Hàn Phi ảnh hưởng sâu sắc đến nhà Tần thống nhất Trung Hoa — và cũng đến tư duy thể chế hiện đại: không ai là không thể thay thế, hệ thống phải tự vận hành.
\textit{(Han Fei's thought profoundly influenced the Qin dynasty that unified China — and also modern institutional thinking: no one is irreplaceable; the system must run itself.)}

\end{truyen}

\begin{baihoc}
Đừng xây tổ chức phụ thuộc vào một cá nhân xuất sắc. Hãy xây hệ thống, quy trình, và văn hóa có thể vận hành tốt ngay cả khi người tài không còn đó.
\textit{(Do not build organizations that depend on one outstanding individual. Build systems, processes, and culture that function well even when the talented person is gone.)}
\end{baihoc}
\ngancach

\section{Hàn Phi Luận Về Lòng Tin Và Quyền Lực}
\begin{truyen}{Hàn Phi Luận Về Lòng Tin Và Quyền Lực}{Han Fei on Trust and Power}
\chuhoa{H}{àn Phi viết: ``Người cai trị khôn ngoan không tin vào lòng tốt của thần tử — ông ta xây dựng hệ thống khiến ngay cả kẻ không tốt cũng hành động đúng.''}
\textit{(Han Fei wrote: ``The wise ruler does not trust in the goodness of his ministers — he builds a system where even the not-good are compelled to act rightly.'')}

Ông phân tích mối quan hệ vua-tôi bằng lăng kính lợi ích: ``Vua muốn tôi hiệu quả. Tôi muốn được thưởng. Khi hai lợi ích này thẳng hàng, mọi việc trôi chảy. Khi chúng mâu thuẫn, tai họa xảy ra.''
\textit{(He analyzed the ruler-minister relationship through the lens of interest: ``The ruler wants ministers to be effective. The minister wants reward. When these two interests align, things flow. When they conflict, disaster follows.'')}

Một vị vua hỏi bề tôi có trung thành không. Hàn Phi cố vấn: ``Đó là câu hỏi sai. Câu đúng là: ta đã tạo ra hệ thống mà lòng trung thành trở nên có lợi hơn phản bội chưa?''
\textit{(A king asked whether his ministers were loyal. Han Fei advised: ``That is the wrong question. The right one is: have I created a system where loyalty is more beneficial than betrayal?'')}

Nghe có vẻ lạnh lùng — nhưng đây là nền tảng của quản trị tổ chức hiện đại: thiết kế cơ chế khuyến khích, không dựa vào lòng tốt tự phát.
\textit{(This sounds cold — but it is the foundation of modern organizational governance: design incentive mechanisms, do not rely on spontaneous goodwill.)}

\end{truyen}

\begin{baihoc}
Khi xây dựng nhóm hay tổ chức, hỏi: ``Hệ thống khen thưởng và trách nhiệm của tôi có đang khuyến khích hành vi đúng không — hay tôi chỉ đang hy vọng vào lòng tốt?'' Cả hai đều cần, nhưng hệ thống phải đến trước.
\textit{(When building a team or organization, ask: ``Are my reward and accountability systems encouraging the right behavior — or am I just hoping for goodwill?'' Both matter, but the system must come first.)}
\end{baihoc}
\ngancach

\section{Mặc Tử Và Chiếc Cung Của Thợ Làm Tên}
\begin{truyen}{Mặc Tử Và Chiếc Cung Của Thợ Làm Tên}{Mozi and the Arrow-Maker's Bow}
\chuhoa{M}{ặc Tử đưa ra hình ảnh ấn tượng để nói về nghề nghiệp và đạo đức:}
\textit{(Mozi offered a striking image to speak about profession and ethics:)}

``Người làm áo giáp muốn người ta dễ bị thương để bán được nhiều áo. Người làm vũ khí muốn chiến tranh xảy ra để bán được dao kiếm. Người thầy thuốc muốn người ta bệnh để có việc làm.'' Nghe có vẻ độc ác — nhưng thực ra chỉ là logic hệ thống.
\textit{(``The armor-maker wants people to be vulnerable so he can sell more armor. The weapon-maker wants war to break out to sell swords. The physician wants people ill to have work.'' This sounds evil — but it is simply system logic.)}

``Vì vậy, hãy chọn nghề nghiệp theo loại thế giới bạn muốn sống trong — không chỉ theo lợi nhuận.''
\textit{(``Therefore, choose your profession according to the kind of world you want to live in — not merely by profit.'')}

Mặc Tử bản thân là thợ thủ công — một người lao động chân tay. Ông không coi triết học là nghề của tầng lớp quý tộc. Ông dạy những người thợ rèn, thợ mộc, dân thường.
\textit{(Mozi himself was a craftsman — a manual laborer. He did not regard philosophy as a profession for the noble class. He taught blacksmiths, carpenters, and common people.)}

Trường phái Mặc Gia có tổ chức chặt chẽ nhất trong các trường phái Trung Hoa cổ đại — với người đứng đầu được gọi là ``Cự Tử'' và quy tắc kỷ luật nghiêm khắc.
\textit{(The Mohist school was the most tightly organized of all ancient Chinese schools — with a leader called ``Juzi'' and strict discipline.)}

\end{truyen}

\begin{baihoc}
Hãy nhìn nghề nghiệp của bạn và hỏi: ``Thế giới tốt hơn hay xấu đi khi tôi giỏi hơn trong nghề này?'' Đó là câu hỏi đạo đức nghề nghiệp thực sự.
\textit{(Look at your profession and ask: ``Is the world better or worse when I become more skilled in this work?'' That is the true question of professional ethics.)}
\end{baihoc}
\ngancach

\section{Hàn Phi Và Cái Chết Của Một Nhà Tư Tưởng}
\begin{truyen}{Hàn Phi Và Cái Chết Của Một Nhà Tư Tưởng}{Han Fei and the Death of a Thinker}
\chuhoa{H}{àn Phi là người nói lắp, không hùng biện được. Nhưng văn chương của ông sắc bén đến mức Hoàng Đế Tần Thủy Hoàng, sau khi đọc, nói: ``Nếu ta được gặp người viết điều này và cùng ông ấy đi dạo, ta chết cũng không tiếc.''}
\textit{(Han Fei had a stutter and could not speak fluently. But his writing was so sharp that Emperor Qin Shi Huang, after reading it, said: ``If I could meet the author of this and walk with him, I would die without regret.'')}

Sau đó Hàn Phi đến Tần. Nhưng Lý Tư — cũng là học trò của Tuân Tử, ganh tị với tài năng của Hàn Phi — vu cáo ông phản quốc.
\textit{(Han Fei then came to Qin. But Li Si — also a student of Xunzi, jealous of Han Fei's talent — accused him of treason.)}

Hàn Phi bị bỏ tù. Trước khi Tần Thủy Hoàng kịp xét xử, Lý Tư gửi thuốc độc vào ngục. Hàn Phi chết trong tù — bởi chính người bạn học cũ của mình.
\textit{(Han Fei was imprisoned. Before Emperor Qin Shi Huang could hold a hearing, Li Si sent poison to the prison. Han Fei died in jail — by the hand of his own former schoolmate.)}

Tư tưởng của ông về sự nguy hiểm của việc cai trị — rằng người cầm quyền luôn bị bao vây bởi những kẻ tư lợi — cuối cùng đã đúng với chính người tạo ra tư tưởng ấy.
\textit{(His thought about the dangers of rule — that those in power are always surrounded by self-interest — proved true for the very person who created it.)}

Nhưng tư tưởng sống sót còn người. Hàn Phi Tử đến nay vẫn là một trong những tác phẩm chính trị sâu sắc nhất mọi thời đại.
\textit{(But thought outlives the person. Han Feizi remains to this day one of the most profound political texts of all time.)}

\end{truyen}

\begin{baihoc}
Những người thách thức trật tự quyền lực thường trả giá cao. Nhưng tư tưởng đúng đắn có sức sống bền dai hơn bất kỳ quyền lực nào.
\textit{(Those who challenge power structures often pay dearly. But sound thought has a longer life than any power.)}
\end{baihoc}
